% Chapter 1

\chapter{INTRODUCTION}\label{cha:intro}

Unmanned aerial vehicles (UAVs) are a practical option for last-mile delivery where
ground transport is slow or impossible: medical supplies for remote clinics, emergency
drug drops after a disaster, and diagnostic samples that should not spend an hour
shaking in the back of a van. Every one of these flights goes from a dispatch hub to a
delivery point, and the air in between holds buildings, hills, restricted zones and
wind. Three things compete when we choose the route: the distance flown, the battery
spent, and the share of the trip spent in risky air. Improve one and another usually
gets worse.

Multi-objective path planning was designed for exactly this. Such a planner carries
distance, energy and risk separately through the search and hands back a Pareto-optimal
front, a set of routes where no route can gain on one objective without losing on
another~\cite{r11}. The front lists the trade-offs on offer. Which one to fly is left
open.

Most systems take a shortcut here. They merge the three objectives into one weighted
sum, give the scalar cost to a plain planner such as A*, and let the operator set the
weights before each flight. If every mission is the same kind, that is fine. Change the
payload and it breaks: emergency medicine wants a very different weighting from a
fragile blood sample, and both differ from a box of relief supplies~\cite{r1}.
The process of re-tuning and re-running the search algorithm between flights consumes
time, and in the event of an emergency, nobody with the experience of the planner is
available.

Our Payload-Aware Two-Tier UAV Route Planning system takes that step away. Tier~1
builds the Pareto front for a delivery region once by using an exact multi objective
graph search over 3D airspace lattice. This lattice is used because its edge costs
depend on terrain, buildings, restricted zones and the way the wind blows. We keep that front until the
graph changes, the front expires or the weather drifts. Tier~2 runs for each mission.
The operator enters four payload tags, a fixed formula turns them into objective
weights, every stored route is scored, and the top one comes back in well under a
millisecond. This is then exported as an ArduPilot mission and flown in a Gazebo and
ArduPilot software-in-the-loop (SITL) simulation.

\section{MOTIVATION}

Three delivery scenarios that come up again and again in medical and emergency drone
work motivate this project, and their routing needs do not agree with each other.

Fragile medical samples such as tissue biopsies and blood specimens must arrive intact,
and an operator will want them kept out of rough air when a smoother corridor exists.
The priority is low flight risk: little time in turbulent or crosswind-exposed corridors
and few sudden altitude changes. We would take a longer route if in case the air is calmer.

Emergency medicine such as insulin or antivenom has a strict time window. The priority
is the shortest route, so that transit time is as low as possible, and risk and energy
come second. If the delivery window is fifteen minutes, the shortest path wins.

Disaster-relief supplies such as water purification tablets and ration packs are
heavier than medical payloads, and they are neither fragile nor urgent. Here the
priority is low battery drain, so that one charge covers as many drops as possible.

A real delivery service sees all three, sometimes within the same hour. A planner with
fixed weights treats them identically, and one that lets the operator change the weights
pushes an optimisation task onto someone dispatching a drone in a hurry. We could not
find a UAV delivery planner that
switches between these routing strategies from a description of the payload, without
an expert in the loop and without running the search again. That is what motivated the
two-tier design.

A second motivation is wind. None of the planners we reviewed resolve wind direction
per edge, yet for a quadcopter cruising at 36\,km/h a 25\,km/h wind is not a detail:
flying into it costs far more energy than flying with it, and crosswind makes the ride
rough.
Section~\ref{sec:res-wind} shows that at a fixed wind speed, the direction alone changes
the energy of the cheapest available route by 27.5\%. A planner blind to direction gets this
wrong on any windy day, and so does any payload-aware choice built on top of it.

\section{BACKGROUND}\label{sec:background}

\subsection{Building a front versus choosing from it}

A routing front can be built in two ways. Exact Graph Search algorithms like
NAMOA*~\cite{r17} keep all non-dominated paths at each node. Conversely, other
algorithms like NSGA-II~\cite{r8} and MOPSO~\cite{r6} perform better on large search
spaces but use random procedures. Typical route selection strategies involve the use of crowding distance
technique~\cite{r8,r3}, an entropy criterion coupled with a weighted sum~\cite{r1}, or
even manual intervention. None of these techniques consider the payload being
delivered by the drone. In the energy model~\cite{r10}, payload is expressed through
mass which impacts edge cost without affecting the relative trade-offs. Our
investigation suggests that no algorithm has considered the payload metadata for
selecting a particular point on the Pareto front.

\subsection{Evidence from real deployments}

Usage of services highlights the instances in which payload is significant, but the
available information is scant. Swiss Post used the Matternet drones for transporting
laboratory samples from one hospital to another and then discontinued its service after
two incidents occurred in 2019. First incident was in January when the drone faced an
electrical fault, which resulted in an interruption of power supply to the GPS system,
and the drone had to land in Lake Zurich urgently. Second incident happened in May when
parachute of the drone got entangled in its frame, resulting in a descent into trees
near some children playing in kindergarten~\cite{r12}. Both were equipment failures, so they say
nothing about route choice. They do show that laboratory samples are routinely flown
over water and populated ground, where the choice of corridor has consequences.

Laboratory studies of medicines and blood samples flown by drone found no damage on
ordinary flights~\cite{r13,r14,r15}, but none of them varied the route
(Section~\ref{sec:lit-papers}), so in this work the fragility tag is a precaution the
operator sets, not a measured degradation effect.

Zipline's blood-delivery network in Rwanda serves hospitals up to about 80\,km from its
distribution centre, and each aircraft follows a predetermined flight plan~\cite{r16}.
Its engineers treat payload weight as a limit on range. That is the energy side of the
problem Tier~2 addresses: a heavier load shortens how far the drone can go, so the best
route depends on what is on board.

\subsection{Regulation}

In India, small civil drones fly under the Drone Rules, 2021~\cite{r21}, which treat
120\,m above ground as the ceiling for flight without special permission. Ground risk is
commonly ranked with the SORA framework~\cite{r22}. Our airspace model uses both.

\section{PROBLEM DEFINITION}\label{sec:problem}

Existing UAV route planners either apply a fixed geometric selection rule over the
Pareto front that ignores the payload, or fold the objectives into one hand-tuned
weighted sum before the search begins, which fixes the trade-off regardless of what is
being carried. When the payload changes, an expert has to re-tune the weights and run the search again,
because no explicit mapping exists from a payload description to the trade-off it
implies. In the field, that expert is not available.

\newpage
\subsection{Formal statement}

We are given:
\begin{enumerate}
  \item an airspace graph $G = (V, E)$ in which every edge carries a three-objective cost
  vector $\mathbf{c}(e) = [d(e),\, \varepsilon(e),\, \rho(e)]$ for flight distance,
  estimated energy and flight risk;
  \item a Pareto front $\mathcal{P}$ of non-dominated routes computed in advance over
  $G$; and
  \item a payload metadata tuple $(F, W, U, P_\mathrm{lvl})$ holding fragility class
  $F \in \{1,\ldots,5\}$, payload weight class $W \in \{1,2,3\}$, delivery urgency
  $U \in \{1,2,3\}$ and priority level $P_\mathrm{lvl} \in \{1,\ldots,5\}$.
\end{enumerate}
The problem is to select the route $p^* \in \mathcal{P}$ that best meets the
requirements implied by the payload metadata, without expert re-tuning and without
running the multi-objective search again.

The inputs are values an operator can assign from domain knowledge (``this vial is
fragile, $F = 5$; the delivery is not urgent, $U = 1$''). Selection must be
deterministic and reproducible, and it must finish within one second of the tags being
entered. Three further conditions follow from the gaps identified in
Chapter~\ref{cha:lit}. The selected route must depend on $(F, W, U, P_\mathrm{lvl})$ and
not on the geometry of the front alone. The front needs to be reusable during multiple
missions before the next refresh instead of constructing it from scratch each time.
Lastly, the conversion from $(F, W, U, P_\mathrm{lvl})$ to $p^*$ needs to be stated in
terms of a mathematical formula; anything else is unacceptable.

\subsection{Weather dependence}

Weather affects the cost of edges. Let $\omega$ be a measurement of wind speed and
direction, gusts, precipitation, and visibility. An edge then costs
$\mathbf{c}(e;\omega) = [d(e),\, \varepsilon(e;\omega),\, \rho(e;\omega)]$ and the
front becomes $\mathcal{P}(\omega)$. The distance will always be constant, while energy
and risk depend on the intersection angle of the wind with each edge. Selection is done
as before. However, the difficulty is that an old front can be used only for as long as
$\omega$ does not deviate too much from the measurements it was built for; therefore,
the front must be updated by time or when the weather deviates. If $\omega$ is outside
the go/no-go limits, then there is no selection and the mission is rescheduled.

\subsection{Inputs, outputs and success criteria}

Table~\ref{tab:ps-io} specifies the inputs, the outputs, and the tests each of them
needs to pass.

\begin{table}[!htbp]
\centering
\caption{Inputs, outputs, constraints and success criteria}\label{tab:ps-io}
\small
\begin{tabular}{|P{3.3cm}|P{11.2cm}|}
\hline
\textbf{Item} & \textbf{Definition}\\
\hline
Region input & Centre point and box size; terrain (DEM), OpenStreetMap buildings, restricted zones, water and land use for that box\\
\hline
Weather input & Snapshot $\omega$ from Open-Meteo: wind at 80\,m, gust, rain, visibility, WMO weather code\\
\hline
Mission input & Depot, destination and payload tuple $(F, W, U, P_\mathrm{lvl})$\\
\hline
Tier~1 output & Stored Pareto front $\mathcal{P}(\omega)$ with graph hash, weather snapshot and expiry time\\
\hline
Tier~2 output & Selected route $p^*$ with waypoints (latitude, longitude, altitude), weights, ranks and point type, or a HOLD advisory\\
\hline
Mission output & QGC WPL~110 \texttt{.waypoints} and QGroundControl \texttt{.plan} files\\
\hline
Hard constraints & Cruise at least 60\,m above ground and 15\,m above roofs; at most 120\,m above ground; no entry into restricted zones; route energy within 80\% of battery capacity\\
\hline
Tier~1 criteria & Hypervolume and IGD at least as good as every baseline on each scenario; no dominated or duplicate routes in the front\\
\hline
Tier~2 criteria & Selection under 1\,s; zero monotonicity violations over all 225 tag combinations; non-extreme picks for mixed payloads\\
\hline
Simulation criteria & Cross-track error, clearance margin and landing offset of the flown route in SITL and Gazebo\\
\hline
\end{tabular}
\end{table}

\subsection{Scope}

This report covers Phase~I which delivers the complete software system and its
simulation. In Phase~II the same system is transferred onto a Pixhawk-equipped F450
quadcopter with a Raspberry Pi~5 companion computer and is tested in outdoor flights. A
Tier~1 front is trusted only for a bounded window. Once the window runs out, or the
weather moves past set thresholds, it is rebuilt; until then Tier~2 alone serves every
mission. Multi-UAV coordination, dodging moving obstacles and real-time replanning are
outside our scope.

\section{ORGANIZATION OF THE REPORT}

There are six chapters, then the references and three appendices.
Chapter~\ref{cha:intro} gives the motivation and the problem. Chapter~\ref{cha:lit}
reviews various near papers, what they leave out, and the objectives that we took from
that. Chapter~\ref{cha:method} explains the architecture and its seven modules including
the algorithms inside them. Chapter~\ref{cha:results} covers datasets, hardware and
software, the experiments run so far and their results. Chapter~\ref{cha:social} looks
at social impact and sustainability, and Chapter~\ref{cha:conclusion} looks at the plan
for the rest of Phase~I and for Phase~II.
